\section{Cactus MPD: additive optimization and exact comparison}
\label{sec:a-cactus}

A cactus separates its cycles topologically, but this does not make all
forms of optimization equally easy. A maximum terminal drop can be
approximated to any prescribed number of bits in polynomial time,
whereas exact comparison already contains Square-Root Sum for one source
and one sink. The positive result rests on the existence of an optimum
with at most two unsaturated nomination coordinates, both in one cycle;
all other block contributions are then constants that can be
approximated separately. We prove this reduction, then give the
algorithm and the exact-arithmetic classification of the
single-source/sink subclass.

\subsection{Setting and results}
\label{subsec:a-cac-setting}
Throughout this section, $G$ is a connected cactus, $s\ne t$ are
prescribed vertices, and the laws are symmetric quadratic laws with
positive rational resistances $\beta_e$. The domain is the nonempty
balanced box $\nomset=\{b:\ell\le b\le u,\ \one^\top b=0\}$ with finite
rational bounds. The intervals may be shifted and may be singletons.
Write $F_{s,t}(b)=\pi_s(b)-\pi_t(b)$ and
$\OPT=\OPT_{F_{s,t}}(\nomset)$. There are no additional flow or potential
bounds. Existence, uniqueness, and attainment follow from
\cref{thm:a-pre-existence,lem:a-pre-attainment}.

\begin{theorem}[Additive optimization on cacti]
\label{thm:a-cac-additive}
For rational $0<\eps\le1$, one can compute a certified rational interval
$[\underline v,\overline v]$ containing $\OPT$, of width at most $\eps$,
and a rational $\widehat b\in\nomset$ satisfying
$F_{s,t}(\widehat b)\ge\OPT-\eps$, in time polynomial in the input bit
length $N$ and $\log(1/\eps)$. As in
\cref{subsec:a-pre-computation}, $N$ includes the encoding of $\eps$.
\end{theorem}

\begin{theorem}[Exact single-source/sink arithmetic]
\label{thm:a-cac-srs}
Let $\nomset_{st}$ be the balanced box with $b_s\in[0,1]$,
$b_t\in[-1,0]$, and $b_v=0$ for $v\notin\{s,t\}$. The exact
threshold problem $\OPT_{F_{s,t}}(\nomset_{st})\ge\gamma$, for rational
$\gamma$, is polynomial-time many-one equivalent to $\SRS$.
Hardness already holds on simple cacti of maximum degree three, all of
whose cycles are triangles, with positive integer resistances.
\end{theorem}

Neither theorem asserts $\NP$-hardness, a polynomial exact algorithm for
multiple-entry cactus MPD, or rational physical flows or potentials.
The additive theorem
addresses the approximation version of the cactus question recorded by
\citet[pp.~9--10]{pfetsch2026-potential-based-flowsan-overview}; the
single-source/sink theorem classifies an exact subclass.

\subsection{Aggregation and electrical sensitivity}
\label{subsec:a-cac-structure}
The \emph{core} is the union of the blocks on the $s$--$t$ path in the
block--vertex incidence tree. Order these blocks from $s$ to $t$, calling
their two path vertices the entrance and exit.

\begin{lemma}[Exact aggregation to the core]
\label{lem:a-cac-aggregation}
Deleting the core edges partitions $V$ into groups $S_v$, one for each
core vertex $v$, with $S_v$ meeting the core precisely at $v$.
Replacing its nominations by
\[
 \bar b_v=\sum_{a\in S_v}b_a,\qquad
 \bar\ell_v=\sum_{a\in S_v}\ell_a,\qquad
 \bar u_v=\sum_{a\in S_v}u_a
\]
preserves both the attainable balanced core nominations and their
terminal potential differences exactly. Every rational feasible core
nomination has a rational disaggregation computable in polynomial bit
time. The same statements hold for the smoothed laws used below.
\end{lemma}
\begin{proof}
By \cref{lem:a-pre-incidence}, any block outside the path belongs to a
subtree attached at a unique path-block vertex. A component after
removing core edges cannot meet two core vertices: the argument in
\cref{lem:a-pre-block}(b) would close a cycle. Connectivity gives an
attachment for every component. Thus the groups partition $V$.
Summing conservation over $S_v$ cancels all its internal flows and
leaves exactly the divergence of the core edges at $v$. Restriction of
the potential equations therefore gives the unique core state for
$\bar b$, also by \cref{lem:a-pre-block}.

Conversely, sums of independent intervals fill their sum interval:
start every coordinate of $S_v$ at its lower bound and
allocate the remainder $\bar b_v-\bar\ell_v$ successively, adding at
most $u_a-\ell_a$ to coordinate $a$. The remainder starts between zero
and the sum of these widths and finishes at zero. Different groups
share no variables, so global balance is preserved. Existence of the
full state and the restriction argument show that its objective is the
prescribed core objective. Rational additions and minima give the
encoding and running-time claim. The argument does not use the
quadratic form, so it also applies to the smoothed laws.
\end{proof}

Potential or flow bounds would in general restrict which disaggregations
are admissible; they are excluded here. Henceforth $G,\ell,u,b,n,m,\nomset$
denote the aggregated core and its data. By \cref{lem:a-cac-aggregation},
$\OPT$ is unchanged. Aggregation has polynomial encoding length and
does not increase the number of vertices or edges.

\begin{lemma}[Smoothed adjoint and its order]
\label{lem:a-cac-adjoint}
For $\rho>0$, set
$\phi_{e,\rho}(x)=\beta_e(x|x|+\rho x)$ and normalize $\pi_t=0$.
The state is continuously differentiable on the space of balanced
nominations. For $F_\rho(b)=\pi_s(b)$ and balanced $db$,
\begin{equation}
 \begin{aligned}
 dF_\rho&=h^\top db,\qquad L_\rho h=e_s-e_t,\quad h_t=0,\\
 L_\rho&=A\operatorname{diag}(1/r_e)A^\top,\qquad
 r_e=\beta_e(2|x_e|+\rho)>0.
 \end{aligned}
 \label{eq:a-cac-adjoint}
\end{equation}
The unit electrical current associated with $h$ traverses the core
blocks in series. The values of $h$ strictly decrease across each
bridge and along both entrance-to-exit branches of each cycle.
The closed block ranges are ordered; consecutive ones meet just at
their common articulation value and nonconsecutive ones are disjoint.
In the unaggregated network, $h$ is constant on every off-core block.
\end{lemma}
\begin{proof}
The law $\phi_{e,\rho}$ is onto, continuously differentiable, and has
strictly positive derivative, so its inverse $\psi_{e,\rho}$ is
continuously differentiable, with derivative $1/r_e$ at the physical
drop. Delete row $t$ of $A$ to obtain $\widetilde A$. In the coordinates
$\widetilde\pi$ and $\widetilde b$ obtained by deleting coordinate $t$,
conservation is
\[
 \widetilde A\psi_\rho(\widetilde A^\top\widetilde\pi)
       -\widetilde b=0.
\]
Its potential Jacobian is
$\widetilde L_\rho=\widetilde A\operatorname{diag}(1/r_e)\widetilde A^\top$.
For a grounded nonzero vector $y$,
$y^\top L_\rho y=\sum_e(y_u-y_v)^2/r_e>0$: equality would make $y$ constant
by connectivity, hence zero by grounding. The implicit function theorem
therefore gives a local continuously differentiable state, and
uniqueness makes these local maps agree on overlaps, hence on the entire
balanced space.
Differentiation gives $\widetilde L_\rho\,d\widetilde\pi=d\widetilde b$.
Symmetry implies
$dF_\rho=\widetilde h^\top d\widetilde b$, where
$\widetilde L_\rho\widetilde h=\widetilde e_s$. Restoring $h_t=0$ gives
\cref{eq:a-cac-adjoint}; the deleted row follows because both sides
sum to zero.

Put $\jmath_e=(h_u-h_v)/r_e$. This is the passive linear state with nomination
$e_s-e_t$. An off-core subtree has zero load away from its attachment.
Summing conservation gives zero divergence there too, so its restriction
has zero nomination and, by uniqueness, zero current and constant
potential. Across each cut separating earlier from later core blocks,
conservation gives total electrical current one. In a bridge $\jmath =1$ in
the entrance-to-exit direction. In a cycle let $R_1,R_2>0$ be the sums
of the electrical resistances along its two branches. Conservation at
interior vertices gives constant branch currents $\jmath_1,\jmath_2$;
equality of the branch drops and their sum give
$R_1\jmath_1=R_2\jmath_2$ and $\jmath_1+\jmath_2=1$, so both are
positive and every edge on either branch has strictly positive drop. Each interior value lies
strictly between the block endpoints, and adding successive positive
block drops proves the range claims.
\end{proof}

The energy/Hessian identity underlying this adjoint is classical:
\citet[Section IV-A, Lemma 1(a)--(b), equations (19) and (21)]{misra2015-maximum-throughput-problem-in-dissipative}
identify potentials as the conjugate-energy gradient and its Hessian
as the inverse grounded weighted Laplacian. Positive smoothing ensures
that all the derivative weights in our application are finite and
strictly positive, including at zero physical flow.

\begin{lemma}[Saturation faces with at most two pivots]
\label{lem:a-cac-faces}
There is a family $\mathcal F$ of $O(n^2)$ closed faces of $\nomset$,
enumerable with rational arithmetic in polynomial time and independent
of resistances, such that one original MPD maximizer lies in their
union. Each face fixes all coordinates at bounds except at most two
pivots. Two pivots lie in one cycle, one strictly inside each of its
entrance-to-exit branches. A sole pivot may also be an articulation or
terminal. Intersecting a face with balance leaves dimension at most one.
\end{lemma}
\begin{proof}
Fix $\rho>0$ and a maximizer $b^*$ of $F_\rho$, which exists by the
continuity just proved and compactness. For every $b'\in\nomset$ the
one-sided derivative along $b^*+\vartheta(b'-b^*)$ is nonpositive. Hence
$b^*$ maximizes $h^\top b$ over the same polytope, and linear-programming
duality supplies a balance multiplier $\mu_0$ with
\begin{equation}
 h_v>\mu_0\ \Longrightarrow\ b_v^*=u_v,\qquad
 h_v<\mu_0\ \Longrightarrow\ b_v^*=\ell_v.
 \label{eq:a-cac-saturation}
\end{equation}
The multiplier also follows directly from exchanges:
if $b_i^*<u_i$ and $b_j^*>\ell_j$, the feasible direction $e_i-e_j$
gives $h_i\le h_j$. Choose $\mu_0$ between the maximum $h_i$ over
the former indices and the minimum $h_j$ over the latter. If either
set is empty, choose $\mu_0$ beyond the other extreme; if both are
empty any value works. These inequalities imply
\cref{eq:a-cac-saturation}. Fixed coordinates belong to neither set,
so this argument holds even when the feasible polytope has smaller
affine dimension or is a singleton.

By \cref{lem:a-cac-adjoint}, a value strictly inside a block range
can equal $h_v$ at at most one interior vertex per cycle branch, and
at no vertex in any other block. Equality with an articulation or
terminal value singles out only that core vertex. A threshold above
$h_s$ forces all lower bounds; one below $h_t$ forces all upper bounds.
These vectors need not be zero and are retained if balanced.

The following enumeration is independent of resistances. Include the
all-lower and all-upper patterns. At each articulation and at $s,t$,
leave that vertex free, put every preceding vertex at its upper bound,
and every following vertex at its lower bound. For each bridge also
include a gap threshold, with its entrance and all preceding vertices
upper and its exit and all following vertices lower. For each cycle,
put the entrance upper and the exit lower, and fix vertices strictly
outside it upper before and lower after. On each ordered list of branch
interior vertices include all gap positions (the prefix upper, the
suffix lower) and all vertex positions (the preceding prefix upper,
the vertex free, and the following suffix lower). Take every pair of
branch choices. A branch with no interior vertices has one gap choice.
This covers a threshold in the open block range, including equality
with one or two interior vertex values. For a cycle of $k_C$ vertices
there are $O(k_C^2)$ choices. Since $\sum_H(|V(H)|-1)=n-1$, their total
is $O(n^2)$. Replace free coordinates with $\ell_v=u_v$ by their fixed
value, discard empty patterns by balance and interval comparisons, and
optionally discard duplicates. Each retained set is a face defined by
bound equalities inside $\nomset$; extraneous feasible patterns are
harmless. The enumeration covers every smoothed maximizer.

It remains to pass to $\rho=0$ without assuming differentiability of
the unsmoothed state. A fixed spanning-tree routing of $b$ is bounded
uniformly on the compact box. The smoothed energy is
\[
 \energy_\rho(x)=\sum_e\beta_e
       \left(\frac{|x_e|^3}{3}+\frac{\rho x_e^2}{2}\right).
\]
For $0<\rho\le1$, comparison with this routing uniformly bounds the
energy minimizers, since the positive cubic terms bound each flow.
If $\rho_i\to0$ and $b_i\to b$, a cluster point $x$ of the minimizing
flows satisfies $Ax=b$. For any competitor $y$ with $Ay=b$, add the
tree routing of $b_i-b$ to obtain competitors $y_i\to y$ for $b_i$.
Passing to the limit in
$\energy_{\rho_i}(x_i)\le\energy_{\rho_i}(y_i)$ gives
$\energy_0(x)\le\energy_0(y)$. Strict convexity identifies $x$
with the original flow, and tree sums of edge drops give convergence
of normalized potentials. This is the continuity argument of
\cref{lem:a-pre-attainment}, with a uniform energy comparison for the
extra smoothing term.

This joint limit also gives uniform convergence $F_\rho\to F_0$ on the compact box:
otherwise a sequence of nominations with errors bounded away from zero
would have a convergent subsequence, contradicting this joint limit and
continuity of $F_0$. Take a convergent subsequence of smoothed maximizers
and, since $\mathcal F$ is finite, a further subsequence in a single
face. Uniform convergence makes the limit an original maximizer, and
closedness keeps it in that face. Neither strict inequalities for a
limiting adjoint nor positive unsmoothed derivatives are needed.
\end{proof}

\subsection{Additive value computation}
\label{subsec:a-cac-algorithm}
On a face with two pivots $v_1,v_2$, let $\bar s$ be minus the sum of the fixed
coordinates. Balance gives $b_{v_1}=z$, $b_{v_2}=\bar s-z$ and
\begin{equation}
 z\in I=[\max\{\ell_{v_1},\bar s-u_{v_2}\},\ \min\{u_{v_1},\bar s-\ell_{v_2}\}].
 \label{eq:a-cac-parameter}
\end{equation}
An empty interval is discarded and a singleton is a fixed nomination.
With one pivot balance fixes it to $\bar s$; with none it tests the fixed
sum. All these operations are exact rational operations.

For a nonconstant face only the cycle containing $v_1,v_2$ is active:
the total nomination of the pivots is fixed, and they belong to the
same component after deleting the edges of any other block, so every
other block has fixed aggregated nominations by \cref{lem:a-pre-block}.
The objective splits as $F_0(z)=D_0+G_C(z)$, where $D_0$ is the sum of the
fixed blocks' entrance-to-exit drops and $G_C$ is the active cycle drop.

\paragraph{Fixed blocks.}
A bridge flow and its drop are rational. In a fixed cycle, routing on
a spanning path and adding a signed circulation gives
$x_e=\xi_e+\chi_e w$, with rational $\xi_e$ and $\chi_e\in\{-1,1\}$.
The loop function $\sum_e\chi_e\beta_ex_e|x_e|$ is strictly increasing
and has a unique root by \cref{lem:a-pre-cycle}. Its rational breakpoints
are $-\xi_e/\chi_e$; on each intervening interval it is a polynomial
of degree at most two. Exact signs at the breakpoints locate the root,
or show it is a breakpoint. A nonempty open piece cannot be identically
zero by strict increase. Linear pieces, double roots, and zero-flow
breakpoint roots are handled by polynomial isolation rather than by
unconditionally dividing by a quadratic leading coefficient. The root
and every block drop are algebraic of degree at most two: on its sign
piece a drop is a rational polynomial of degree at most two in that
root. Rational root isolation and refinement
\citep[Chapter 10]{basu2006-algorithms-in-real-algebraic-geometry}
approximate each drop in polynomial bit time; once the root is
bracketed, interval evaluation of the drop converges at a rate bounded
by its polynomial derivative on the bracket. All these bounds have
polynomial bit length. The sum $D_0$ is never isolated as one algebraic
number.

\paragraph{Uniform bounds.}
Set
\begin{equation}
 B=\sum_v\max\{|\ell_v|,|u_v|\},\qquad
 \Vmax=B^2\sum_e\beta_e.
 \label{eq:a-cac-bounds}
\end{equation}
Every physical flow, also for $\rho>0$, satisfies $|x_e|\le B$.
To see this, orient each nonzero edge along its actual flow. The
potential strictly decreases along it, so there are no directed cycles. Add a super
source supplying each positive nomination and a super sink collecting
each negative nomination. Conservation permits repeatedly extracting
a positive source-to-sink path, subtracting the minimum remaining flow
on it, and deleting an exhausted edge. There is no residual circulation
because the orientation is acyclic. The path weights sum to total
injection, at most $B$, and each edge flow is a sub-sum of these weights.
Aggregated block nominations are sums over a partition, so their
$\ell_1$ norm is at most $B$, and so is the magnitude of their
spanning-path routing. Subtracting it from the physical cycle
flow proves $|w|\le2B$. Finally any simple path drop, including an active
block drop or the full objective, has absolute value at most $\Vmax$.
If $B=0$, the only core nomination and all core flows are zero, and the
algorithm returns $[0,0]$ and a disaggregation immediately.

\paragraph{Active cycle.}
For a $k_C$-edge active cycle, rational tree routing now gives
\begin{equation}
 x_e(z,w)=\xi_e+\omega_e z+\chi_ew.
 \label{eq:a-cac-active-flow}
\end{equation}
The $k_C$ flow-zero lines form an arrangement with $O(k_C^2)$ cells. One
can enumerate it by inserting lines successively: a new line has at
most $k_C-1$ distinct intersection points and hence at most $k_C$ segments
splitting existing cells. Intersections, ordering along lines, and
side tests are rational computations; coincident lines can be merged.
Use closures of the two-dimensional cells, retaining each original
line's sign. These closures cover the whole plane, including all
lower-dimensional intersections. Intersect with $z\in I$ and
$-2B\le w\le2B$; this also covers singleton $I$ and zero flows.

On a cell let $\epsilon_e\in\{-1,1\}$ specify the weak sign of $x_e$.
The physical loop condition and the active drop along one fixed
entrance-to-exit branch are respectively
\[
 \begin{aligned}
 \Phi_C(z,w)&=\sum_e\chi_e\beta_e\epsilon_ex_e(z,w)^2=0,\\
 G_C(z,w)&=\sum_{e\text{ on branch}}\delta_e\beta_e\epsilon_ex_e(z,w)^2.
 \end{aligned}
\]
where $\delta_e$ is the branch traversal sign. They have degree at most
two. The cell inequalities $\epsilon_ex_e\ge0$ and the compact box are
linear. Conservation is already built into
\cref{eq:a-cac-active-flow}; the loop condition is sufficient for a
potential by \cref{lem:a-pre-cycle}. At a zero flow both polynomial
pieces agree, so closed cells add no false states.

For rational $\gamma_*$, feasibility with the further inequality $G_C(z,w)\ge\gamma_*$
is an existential semialgebraic problem in \emph{two} real variables,
with $O(k_C)$ constraints of degree at most two per cell. Clear positive
denominators first. Aggregation, routing, squaring affine forms, and
summing at most $m$ terms use sums and fixed powers of input rationals;
a common denominator gives coefficient bit length $O(N+\iota)$ after
$\iota$ threshold bisections. The displayed estimates use the safe
upper bound $O(N^4+\iota)$. By the Renegar bound
in \cref{subsec:a-pre-computation}, a cell decision costs at most
\[
 O\bigl(k_C^{a_{\rm dec}}(N^4+\iota+1)^2\bigr)
\]
bit operations, for an absolute constant exponent $a_{\rm dec}$ (the exponent
for two variables and degree two). Formula evaluation is included.
Alternatively \cref{eq:a-pre-qe-cost} with two quantified variables and
one unused free variable gives polynomial-count exponent six and an
absolute constant degree factor; its coefficient-bit bound is polynomial
as well. Neither exponent depends on cactus size.

Every $z\in I$ has a feasible point in the union of the cells. Binary
search on $[-\Vmax,\Vmax]$ using these exact tests gives an interval for
$\max_I G_C$ of width $\eps/8$. Keep a feasible lower threshold and a
cell witnessing it. The initial lower endpoint $-\Vmax$ is feasible; the
upper endpoint need only bound the optimum, so equality at $\Vmax$ causes
no problem. At most $O(1+\log^+(\Vmax/\eps))$ iterations suffice, where
$\log^+ a=\max\{0,\log_2 a\}$.
Approximate each fixed block drop by a rational interval of width at
most $\eps/(8n)$. There are at most $n-1$ blocks, so the resulting
interval for $D_0$ has width at most $\eps/8$. Adding gives a face
interval $[\underline v_J,\overline v_J]$ of width at most $\eps/4$;
fixed nominations use only the fixed-block procedure. Then
\begin{equation}
 [\underline v,\overline v]
       =[\max_{J\in\mathcal F}\underline v_J,\ \max_{J\in\mathcal F}\overline v_J]
 \label{eq:a-cac-interval}
\end{equation}
contains the global optimum by \cref{lem:a-cac-faces}. Its width is
at most $\eps/4$: choose a face supplying the largest upper endpoint,
and note that the largest lower endpoint is at least its lower endpoint.

\begin{algorithm}[Cactus additive MPD]
\label{alg:a-cac-additive}
Input the data of \cref{thm:a-cac-additive}.
\begin{enumerate}
\item Find the block path, aggregate intervals, and retain the groups
for disaggregation. Enumerate the saturation faces of
\cref{lem:a-cac-faces}, discard empty ones, and parameterize them by
\cref{eq:a-cac-parameter}. Handle $B=0$ directly.
\item For each face, isolate and refine fixed-block drops to total
interval width $\eps/8$. If the face varies, enumerate its active
cycle's sign cells and bisect the active optimum to width $\eps/8$,
retaining a feasible lower threshold and cell.
\item Form the face intervals and return the value interval
\cref{eq:a-cac-interval}. Select a face with greatest lower endpoint.
\item On that face recover a feasible algebraic parameter at the
retained lower threshold, round it within $\eps/(8C)$ as in the proof of
\cref{thm:a-cac-additive}, recompute the other pivot by balance, and
greedily disaggregate. A fixed face already has a rational nomination.
Return this rational nomination together with the value interval.
\end{enumerate}
\end{algorithm}

\subsection{Rational nomination recovery}
\label{subsec:a-cac-recovery}
\begin{lemma}[Lipschitz estimate]
\label{lem:a-cac-lipschitz}
For any simple $s$--$t$ core path $P_{st}$ and balanced $b,b'\in\nomset$,
\[
 |F_0(b)-F_0(b')|\le C\norm{b-b'}_1,
 \qquad C=2B\sum_{e\in P_{st}}\beta_e.
\]
This estimate also holds on any connected graph with the same data
assumptions, using a simple terminal path and its nomination bound $B$.
\end{lemma}
\begin{proof}
The flow bound in \cref{eq:a-cac-bounds} gives derivative resistances
$r_e\le\beta_e(2B+\rho)$. We prove the two electrical facts needed to
bound the adjoint. At a vertex other than $s,t$, $L_\rho h=0$ expresses $h_v$
as a weighted average of its neighbors with positive weights. A global
maximum at such a vertex propagates to its neighbors until it reaches
a terminal. The sink cannot be a maximum since $(L_\rho h)_t=-1$, and the
source cannot be a minimum since $(L_\rho h)_s=1$. Applying the same argument
to the minimum shows $0=h_t\le h_v\le h_s$ everywhere.

For $\jmath=\operatorname{diag}(1/r_e)A^\top h$, we have
$A\jmath=e_s-e_t$ and $\sum_er_e\jmath_e^2=h^\top A\jmath=h^\top(e_s-e_t)=h_s$. For any other unit
terminal flow $y$, write $y=\jmath+\Delta x$, with $A\Delta x=0$. Then
\[
 \sum_er_ey_e^2=\sum_er_e\jmath_e^2+\sum_er_e(\Delta x_e)^2,
\]
because the cross term is $2h^\top A\Delta x=0$. Testing with the signed unit
flow on $P_{st}$ proves the effective-resistance/path comparison
\[
 0\le h_v\le h_s\le\sum_{e\in P_{st}}r_e
                   \le(2B+\rho)\sum_{e\in P_{st}}\beta_e.
\]
Integrate $h^\top(b'-b)$ along the balanced feasible segment from $b$
to $b'$, using \cref{lem:a-cac-adjoint}. The absolute value is bounded
by the last upper bound times $\norm{b'-b}_1$. Uniform convergence of
the smoothed objectives, proved in \cref{lem:a-cac-faces} by an argument
valid on any connected graph, permits $\rho\downarrow0$.
\end{proof}

\begin{proof}[Proof of \cref{thm:a-cac-additive}]
It remains to justify witness extraction, error budgets, and bit costs
in \cref{alg:a-cac-additive}. For a selected nonconstant face, take the
cell and feasible lower threshold $\gamma_-$ from its bisection. Eliminate
$w$ from the cell constraints, $\Phi_C=0$, and $G_C\ge \gamma_-$, leaving a
quantifier-free formula in $z$. The bound
\cref{eq:a-pre-qe-cost} has one quantified and one free variable here,
so its polynomial-count exponent is four, with constant output degree
and polynomial coefficient bits. Isolate the roots of these univariate
output polynomials and test their signs at the roots and in the
intervening intervals, using
\citet[Chapter 10]{basu2006-algorithms-in-real-algebraic-geometry}.
This is the one-variable case of the fixed-dimension sample-point
primitive in \cref{prim:a-pre-sample-points}.
This yields an algebraic feasible $z^*$ of polynomial encoding length;
a feasible interval may instead supply a rational point. Identically
zero polynomials and singleton feasible sets are covered by sign tests.
A corresponding $w$ exists by the projected formula and is the unique
physical circulation, so neither the algebraic constant sum nor $w$
needs to be compared or output. The resulting objective is at least the
selected face's certified lower endpoint: its active drop is at least
$\gamma_-$, and its constant contribution is at least the lower endpoint of
the interval for $D_0$.

If $C>0$, refine $z^*$ and clip a rational approximation to its rational
interval $I$, obtaining $|\widehat z-z^*|\le\eps/(8C)$. Clipping cannot
increase distance to a point in $I$. Set the other pivot to $\bar s-\widehat z$.
The nomination changes in $\ell_1$ norm by at most $\eps/(4C)$, so
\cref{lem:a-cac-lipschitz} limits the loss to $\eps/4$. All interval
endpoints, fixed coordinates, and singleton cases remain exactly
feasible. Since $s\ne t$ and resistances are positive, $C=0$ means
$B=0$, already handled. Greedy disaggregation from
\cref{lem:a-cac-aggregation} gives rational original loads satisfying
all original bounds and balance exactly, with the same objective.

Selecting the greatest lower endpoint loses at most $\eps/4$ relative
to the global optimum: if face $J$ contains a global maximizer, then
$\OPT\le \overline v_J\le \underline v_J+\eps/4\le \underline v_{\rm selected}+\eps/4$.
The extracted parameter attains at least $\underline v_{\rm selected}$,
and rounding loses at most $\eps/4$. Thus the nomination loss is at
most $\eps/2\le\eps$, while
\cref{eq:a-cac-interval} is a certified value interval of width at
most $\eps/4\le\eps$. This is precisely the value and exactly
admissible rational-scenario contract of \cref{def:a-pre-problems};
no rational state output is asserted.

Let $N_\eps=N+\lceil\log_2(1/\eps)\rceil+1$. The rational
bounds $B,\Vmax,C$, interval endpoints, and isolated quadratic roots have
polynomial bit length; value bisection uses $O(N_\eps)$ iterations and
rounding uses polynomially many bits. There are $O(n^2)$ faces and
$O(n^2)$ cells per face. With the absolute constant $a_{\rm dec}$ above, their
decision cost is bounded by
\[
 O\bigl(N^{a_{\rm dec}+4}(N^4+N_\eps)^2N_\eps\bigr).
\]
All other rational work, fixed quadratic isolation/refinement, and the
one-variable projection and recovery have cost $O((N+N_\eps)^{a_{\rm rec}})$ for
another absolute constant $a_{\rm rec}$, by the stated fixed-dimension bounds
and univariate isolation. These exponents are independent of the
number of blocks, vertices, and accuracy bits. The sum is a polynomial
in $N$ and $\log(1/\eps)$, as required.
\end{proof}

\subsection{Exact arithmetic and the existential-real upper bound}
\label{subsec:a-cac-arithmetic}
\begin{lemma}[Parallel branches and the radical gadget]
\label{lem:a-cac-radical}
Two internally unloaded branches of total quadratic resistances
$\varrho_1,\varrho_2>0$ have effective quadratic resistance
\begin{equation}
 \begin{aligned}
 R(\varrho_1,\varrho_2)
 &=\frac{\varrho_1\varrho_2}{(\sqrt{\varrho_1}+\sqrt{\varrho_2})^2}\\
 &=\begin{cases}
 \displaystyle\frac{\varrho_1\varrho_2(\varrho_1+\varrho_2)}{(\varrho_1-\varrho_2)^2}
       -\frac{2\varrho_1\varrho_2}{(\varrho_1-\varrho_2)^2}\sqrt{\varrho_1\varrho_2},&\varrho_1\ne \varrho_2,\\[5pt]
 \varrho_1/4,&\varrho_1=\varrho_2.
 \end{cases}
 \end{aligned}
 \label{eq:a-cac-parallel}
\end{equation}
For integer $a\ge2$, the choice $\varrho_1=(a-1)^2$, $\varrho_2=(a-1)^2/a$
has unit-flow drop $a+1-2\sqrt a>0$.
\end{lemma}
\begin{proof}
Internal conservation makes each branch flow constant. For total
$\zeta>0$ both flows have the sign of their common potential drop, so both
are positive. Equal drops and conservation give
\[
 x_1=\frac{\zeta\sqrt{\varrho_2}}{\sqrt{\varrho_1}+\sqrt{\varrho_2}},\qquad
 x_2=\frac{\zeta\sqrt{\varrho_1}}{\sqrt{\varrho_1}+\sqrt{\varrho_2}}.
\]
Their common drop is $\zeta^2R(\varrho_1,\varrho_2)$. At $\zeta=0$ uniqueness gives zero;
symmetry treats negative through-flow if needed. For $\varrho_1\ne \varrho_2$,
multiply $1/(\varrho_1+\varrho_2+2\sqrt{\varrho_1\varrho_2})$ by
$(\varrho_1+\varrho_2-2\sqrt{\varrho_1\varrho_2})/(\varrho_1+\varrho_2-2\sqrt{\varrho_1\varrho_2})$; the denominator becomes
$(\varrho_1-\varrho_2)^2$. If $\varrho_1=\varrho_2$, the equal split gives $\varrho_1/4$ directly.
For the gadget $x_1=1/(1+\sqrt a)$ at unit flow, whence
$(a-1)^2x_1^2=(\sqrt a-1)^2=a+1-2\sqrt a>0$.
\end{proof}

The parallel formula is classical; it is the quadratic specialization
of the homogeneous reduction in
\citet[Lemma 4.2]{gro2017-algorithmic-results-for-potential-based}.
Its use for the following arithmetic reductions is what matters here.

\begin{proof}[Proof of \cref{thm:a-cac-srs}]
Write $\Delta_{st}$ for the unit-flow terminal drop on the network under
consideration. For $\nomset_{st}$, balance writes every nomination
as $\zeta(e_s-e_t)$ with $0\le \zeta\le1$. Homogeneity
(\cref{lem:a-pre-homogeneity}) gives a drop equal to $\zeta^2$ times the
unit-flow drop. The latter is positive since
\[
 \pi_s-\pi_t=\sum_e\beta_e|x_e|^3>0
\]
at unit nomination by \cref{eq:a-pre-dissipation}. Hence the maximum
is attained at $\zeta=1$, and
$\OPT_{F_{s,t}}(\nomset_{st})=\Delta_{st}$.

\emph{Reduction from $\SRS$.}
Start with $\sum_i\sqrt{a_i}\le K$ in the weak convention of
\cref{subsec:a-pre-computation}. Remove any zero radicands if allowed;
remove each radicand one and subtract one from $K$. If no terms
remain, output a fixed yes/no MPD instance according to $0\le K$:
a resistance-one edge with threshold one/two. Otherwise all remaining
$a_i\ge2$; denote their number by $\mu\ge1$. Make $\mu$ disjoint triangles.
Triangle $i$ has direct-edge resistance $(a_i-1)^2$ and two
alternate-path edges each of resistance $(a_i-1)^2/(2a_i)$.
Join successive exits and entrances by unit-resistance bridges, with
$s$ the first entrance and $t$ the last exit. This is a simple cactus
with $3\mu$ vertices and $4\mu-1$ edges. Every triangle terminal has two
triangle edges and at most one bridge, so maximum degree is three.
Every block carries unit through-flow at $\zeta=1$, giving
\[
 \Delta_{st}=C_0-2\sum_{i=1}^\mu\sqrt{a_i},\qquad
 C_0=\sum_{i=1}^\mu(a_i+1)+(\mu-1).
\]
Set $\gamma=C_0-2K$. Then $\Delta_{st}\ge\gamma$ if and only if the
preprocessed SRS inequality holds, including equality. If the
preprocessed $K\le0$, it is a no instance, since a positive radical
remains, and one may output the fixed no instance. If $\gamma\le0$,
the constructed instance already answers yes, and one may instead
output the fixed yes instance: the constructed unit-flow drop
is strictly positive, so $C_0>2\sum_i\sqrt{a_i}$ and $K\ge C_0/2$ imply yes.
This handles nonpositive thresholds without changing the weak convention.

In the remaining construction multiply every resistance and $\gamma$
by $\Lambda=2\prod_i a_i$. Flows are unchanged, since multiplying the
potential vector by $\Lambda$ gives precisely the scaled laws; uniqueness
then applies. All resistances become positive integers and the threshold
becomes an integer. The bit length of $\Lambda$ is
$O(1+\sum_i\log a_i)$, and each of the linearly many output numbers
has polynomial bit length. Only rational arithmetic, integer
multiplication, and exact division are required.

\emph{Reduction to $\SRS$.}
On an arbitrary single-source/sink cactus, every off-path block has
zero aggregated nomination and therefore zero flow by
\cref{lem:a-pre-block,thm:a-pre-existence}. The path blocks form a
series chain. A bridge contributes its resistance at unit flow; a
cycle contributes \cref{eq:a-cac-parallel}, where $\varrho_1,\varrho_2$ are sums of
its branch resistances. Equal branches contribute the rational $\varrho_1/4$.
For each unequal pair absorb the positive radical coefficient into
its radicand, obtaining
\begin{equation}
 \Delta_{st}=Q-\sum_{j=1}^{\mu'}\sqrt{r_j},\qquad
 r_j=\frac{4\varrho_{1,j}^3\varrho_{2,j}^3}{(\varrho_{1,j}-\varrho_{2,j})^4}>0,
 \label{eq:a-cac-radical-sum}
\end{equation}
where $Q>0$ is rational. If $\mu'=0$, decide the rational comparison and
output the fixed SRS yes instance $(a_1,K)=(1,1)$ or no instance
$(4,1)$. Otherwise put $T=Q-\gamma$. If $T\le0$, output that fixed
no instance, because every radical is positive. This includes $T=0$
and respects the positive-integer threshold convention for SRS.
For $T>0$, write $r_j=\mathfrak a_j/\mathfrak b_j$ and $T=\mathfrak p/\mathfrak q$ with positive integer
denominators and define
\[
 \mathcal D=\mathfrak q\prod_{j=1}^{\mu'} \mathfrak b_j,\qquad
 a'_j=\mathcal D^2\mathfrak a_j/\mathfrak b_j\in\Z_{>0},\qquad K'=\mathcal D T\in\Z_{>0}.
\]
Multiplication by positive $\mathcal D$ gives the exact equivalence
$\Delta_{st}\ge\gamma\ \Longleftrightarrow\ \sum_j\sqrt{a'_j}\le K'$.
Branch sums, their differences, and the fixed powers in
\cref{eq:a-cac-radical-sum} have polynomial rational encoding length.
Denominator products have bit length at most the sum of factor bit
lengths, so $Q,\mathcal D,a'_j,K'$ and the complete output also have polynomial
encoding length. A small nonzero $\varrho_1-\varrho_2$ changes magnitudes but not this
bound. Testing $\varrho_1=\varrho_2$ is exact rational comparison. No factorization or
square-root evaluation occurs in either encoder.
\end{proof}

\begin{proposition}[General exact MPD is existential-real]
\label{prop:a-cac-etr}
For any connected graph with positive rational symmetric quadratic
resistances, rational $\gamma$, and a balanced rational nomination box,
the exact threshold problem
$\exists b\in\nomset:\pi_s(b)-\pi_t(b)\ge\gamma$ is in $\ETR$, and
therefore in $\PSPACE$.
\end{proposition}
\begin{proof}
Introduce real variables $b_v,\pi_v$ and $p_e,n_e$ and impose
\begin{align*}
 &\ell_v\le b_v\le u_v\quad(v\in V),\qquad
 \one^\top b=0,\qquad A(p-n)=b,\qquad\pi_t=0,\\
 &p_e\ge0,\quad n_e\ge0,\quad p_en_e=0,\quad
 \pi_u-\pi_v=\beta_e(p_e^2-n_e^2)
       \quad(e=(u,v)),\\
 &\pi_s-\pi_t\ge\gamma.
\end{align*}
For nonnegative complementary $p_e,n_e$, at least one is zero.
If $p_e>0$ then $x_e=p_e-n_e=p_e$ and $x_e|x_e|=p_e^2$;
if $n_e>0$ then $x_e=-n_e$ and $x_e|x_e|=-n_e^2$; when both
vanish both sides are zero. Conversely every signed $x_e$ has the
unique such split $p_e=\max\{x_e,0\}$, $n_e=\max\{-x_e,0\}$.
Thus the constraints encode exactly the passive equations and the
threshold, and existence and uniqueness from \cref{thm:a-pre-existence}
identify their potential difference with $F_{s,t}(b)$.
There are $2n+2m$ variables and $O(n+m)$ constraints of degree at most
two and polynomial coefficient bit length; clearing positive rational
denominators preserves polynomial size. Existentially quantifying all
variables gives the reduction to $\ETR$.
The real-algebraic procedures in \cref{subsec:a-pre-computation} give
decidability; for the stronger space bound with unbounded dimension
we use the polynomial-space existential-real decision theorem of
\citet[report version, Theorem 3.3]{canny1988-some-algebraic-and-geometric-computations}.
\end{proof}

Consequently exact cactus MPD threshold comparison is $\SRS$-hard,
already for one source and one sink, and is in $\ETR\subseteq\PSPACE$.
Whether the problem with multiple entries is polynomial-time equivalent
to $\SRS$ remains open here.

\subsection{Precision and earlier algorithms}
\label{subsec:a-cac-discussion}
Additive precision does not settle an equality-sensitive threshold:
when the certified interval straddles $\gamma$, narrowing it need not
terminate with an answer, particularly if $\gamma=\OPT$.
\Cref{thm:a-cac-additive} gives no bound on the number of accuracy bits
needed for a general exact comparison.

\citet[Section 6]{labbe2020-bookings-in-the-european-gas}, by Labb\'e,
Plein, and Schmidt, leaves booking feasibility on a single cycle with
nonlinear potentials open.
The later four-author paper of Labb\'e, Plein, Schmidt, and Th\"urauf
\citep[Section 6, Theorems 6.4--6.5 and Corollary 6.6
(numbering of the November 2020 manuscript)]{labbe2021-deciding-feasibility-on-a-cycle}
gives polynomial Turing-time booking validation on a passive cycle
under rational-data Assumption 6.1. Its structural configurations reduce
the decision systems to at most nine variables and 42 constraints.
Thus both structural nomination restrictions and fixed-dimensional
real algebraic computation have direct precedents. With arbitrarily
many cactus cycles, even fixed single-source/sink nominations can
accumulate an unbounded number of independent radicals. Exact local
solutions and a compact sum expression do not themselves supply
polynomial exact threshold comparison.

The series-parallel reductions of
\citet[Lemma 4.2 and Theorem 4.3]{gro2017-algorithmic-results-for-potential-based}
(later published in \emph{Networks}, 2019) construct an equivalent arc with
linearly many reduction applications. Their Section 2.1 explicitly
uses the Turing model and discusses root approximation with a fixed
error, so their result is not merely a real-arithmetic statement.
The stronger zero-gap threshold guarantee needs a separate argument,
which \cref{thm:a-cac-srs} identifies with SRS on the stated subclass.
\citet[Lemma 4.20 and Section 5]{thurauf2022-deciding-the-feasibility-of-a}
proves general-network MPD hardness, also with predetermined flow
directions, and proposes both cacti and approximate MPD as further
directions; the overview retains the nonlinear cactus question
\citep[pp.~9--10]{pfetsch2026-potential-based-flowsan-overview}.
The results here give polynomial additive optimization for balanced
shifted boxes and an exact arithmetic classification for a single
source and sink, while leaving the multiple-entry exact boundary open.
